% =====================================================================
% RELATÓRIO 10 — POO com JavaScript
% Capítulo para ser incorporado à Apostila S01 (OO + PF), após o Relatório 9.
% =====================================================================

% ---------------------------------------------------------------------
% Definições usadas neste capítulo (podem ser movidas para o preâmbulo).
% - JavaScript: o listings não tem essa linguagem, então ela é definida
%   aqui (palavras-chave, comentários e strings). As strings entre aspas
%   duplas ficaram de fora de propósito: os códigos deste capítulo só
%   usam aspas simples e crases, e depois do begin{document} o babel em
%   português transforma as aspas duplas em caractere ativo, o que pode
%   quebrar a definição.
% - javascript: estilo que usa a linguagem acima e corrige a acentuação
%   dentro dos códigos (mesma solução dos estilos csharp, python e ruby).
% - azulRelatorio e exercicio: já são criados no Relatório 7. As linhas
%   abaixo só os recriam se este capítulo for compilado sem ele.
% ---------------------------------------------------------------------
\lstdefinelanguage{JavaScript}{
    keywords={class, constructor, extends, super, this, new, return, function, const, let,
        var, if, else, for, of, in, while, static, get, set, true, false, null, undefined,
        typeof, instanceof, throw, try, catch},
    sensitive=true,
    comment=[l]{//},
    morecomment=[s]{/*}{*/},
    morestring=[b]',
    morestring=[b]`
}

\lstdefinestyle{javascript}{
    language=JavaScript,
    literate=
        {á}{{\'a}}1 {à}{{\`a}}1 {ã}{{\~a}}1 {â}{{\^a}}1
        {é}{{\'e}}1 {ê}{{\^e}}1 {è}{{\`e}}1 {í}{{\'\i}}1
        {ó}{{\'o}}1 {õ}{{\~o}}1 {ô}{{\^o}}1 {ú}{{\'u}}1
        {ç}{{\c{c}}}1 {Á}{{\'A}}1 {É}{{\'E}}1 {Ç}{{\c{C}}}1
        {Ã}{{\~A}}1 {Â}{{\^A}}1 {Ê}{{\^E}}1 {Õ}{{\~O}}1
        {—}{{\textemdash}}1 {Í}{{\'I}}1 {Ó}{{\'O}}1 {Ú}{{\'U}}1
}

\definecolor{azulRelatorio}{rgb}{0.1,0.3,0.8}

\ifcsname exercicio\endcsname\else
\newtcolorbox{exercicio}[1]{
    colback=white,
    colframe=azulRelatorio,
    coltitle=white,
    fonttitle=\normalfont,
    boxrule=0.8pt,
    arc=1.5mm,
    title={#1}
}
\fi

\chapter{Relatório 10: Programação Orientada a Objetos em JavaScript}

\section{Fundamentos Teóricos e Pilares da POO}

No Relatório 9, o Ruby mostrou uma linguagem em que tudo é objeto e em que os blocos terminam
com \texttt{end}. O \textbf{JavaScript} nasceu em outro contexto: foi criado para rodar dentro do
navegador e reagir ao que o usuário faz, e hoje também roda fora dele, em servidores, com o
\textbf{Node.js}. Assim como o Python e o Ruby, ele tem \textbf{tipagem dinâmica}: as variáveis
não declaram tipo. Já a sintaxe lembra o C\# e o C++: os blocos usam chaves e as instruções
terminam com ponto e vírgula. A palavra \texttt{class} foi acrescentada à linguagem em 2015 e
deixa o código parecido com o das linguagens anteriores. Os objetos são criados com
\texttt{new}, os atributos são acessados sempre com \texttt{this} e os atributos privados
começam com \texttt{\#}. Assim como nas linguagens anteriores, a memória é liberada
automaticamente pelo \textbf{Garbage Collector}.

A tabela abaixo relaciona as estruturas já vistas em Ruby com as equivalentes em JavaScript:

\begin{center}
\begin{tabular}{|l|l|l|}
\hline
\textbf{Conceito} & \textbf{Ruby (Relatório 9)} & \textbf{JavaScript (Relatório 10)} \\
\hline
Construtor & \texttt{def initialize(nome)} & \texttt{constructor(nome)} \\
Atributo do objeto & \texttt{@nome} & \texttt{this.nome} \\
Instanciação & \texttt{Item.new(nome)} & \texttt{new Item(nome)} \\
Atributo privado & toda variável \texttt{@} & prefixo \texttt{\#} (declarado) \\
Getter & \texttt{attr\_reader :nome} & \texttt{get nome() \{ ... \}} \\
Herança & \texttt{class Arma < Item} & \texttt{class Arma extends Item} \\
Construtor do pai & \texttt{super(nome)} & \texttt{super(nome)} \\
Método do pai & \texttt{super} & \texttt{super.metodo()} \\
Interface & \texttt{module} + \texttt{include} & \texttt{throw new Error()} na base \\
Verificação de classe & \texttt{x.is\_a?(Item)} & \texttt{x instanceof Item} \\
Inserção na coleção & \texttt{lista <{}< item} & \texttt{lista.push(item)} \\
Repetição & \texttt{lista.each do |x| ... end} & \texttt{lista.forEach(x => \{ ... \})} \\
Filtro & \texttt{lista.select \{ |x| ... \}} & \texttt{lista.filter(x => ...)} \\
Fim de bloco & \texttt{end} & \texttt{\}} \\
\hline
\end{tabular}
\end{center}

Os quatro pilares continuam valendo, agora ilustrados com uma banda de rock que se prepara para
um show:

\begin{itemize}
    \item \textbf{Abstração}: Cada objeto guarda apenas o que importa para o show. Um
    \texttt{Musico} precisa de nome, função e experiência.

    \item \textbf{Encapsulamento}: A experiência do músico fica no atributo privado
    \texttt{\#experiencia}, que só pode ser lido pelo getter \texttt{experiencia}. Na
    \texttt{Banda}, os membros e o repertório também são privados.

    \item \textbf{Herança}: \texttt{Vocalista} reaproveita a estrutura de \texttt{Musico}, e
    \texttt{Guitarra} e \texttt{Bateria} reaproveitam a de \texttt{Instrumento} (relação
    ``\textbf{É UM}'': um Vocalista é um Músico).

    \item \textbf{Polimorfismo}: No show, a banda chama \texttt{executarAcao()} de cada membro e
    \texttt{tocar()} de cada instrumento, e cada objeto executa a sua própria versão do método.
\end{itemize}

\subsection{Interfaces Simuladas e \textit{Duck Typing}}

O JavaScript não tem a palavra-chave \texttt{interface} nem classes abstratas nativas. O
contrato é simulado com uma classe base cujo método apenas \textbf{lança um erro} com
\texttt{throw new Error(...)}. As classes filhas sobrescrevem esse método; se alguma não
sobrescrever, o erro aparece quando o método for chamado. No código de referência, a classe
\texttt{Instrumento} exige que \texttt{Guitarra} e \texttt{Bateria} implementem
\texttt{tocar()}.

Assim como no Ruby, e diferente do \texttt{ABC} do Python, nada é conferido antecipadamente: a
classe base pode ser instanciada normalmente, e o erro só aparece em tempo de execução. Também
como no Ruby, o que importa é se o objeto \textbf{sabe responder} ao método chamado (o
\textit{duck typing}). No código de referência, um objeto criado na hora, que não pertence a
nenhuma classe, funciona como instrumento só porque tem o método \texttt{tocar()}.

\subsection{Relacionamentos entre Classes}

A composição e a agregação, vistas nos relatórios anteriores, continuam valendo (relação
``\textbf{TEM UM}''):

\begin{itemize}
    \item \textbf{Composição (relação forte)}: A parte é criada \textbf{dentro} do objeto dono e
    ``morre'' junto com ele. Exemplo: o repertório da banda é criado dentro do
    \texttt{constructor} da própria \texttt{Banda}.

    \item \textbf{Agregação (relação fraca)}: A parte é criada \textbf{fora} e o dono apenas
    guarda uma referência a ela. Exemplo: os músicos são criados antes e só depois entregues à
    \texttt{Banda}. O mesmo vale para os instrumentos, criados no programa principal e depois
    entregues aos músicos.
\end{itemize}

\section{Documentação Sintática das Estruturas de POO}

Esta seção detalha todas as palavras-chave e construções sintáticas utilizadas na implementação
de POO em JavaScript. Os exercícios deste relatório devem ser resolvidos \textbf{apenas} com as
estruturas documentadas aqui.

\subsection{Estrutura do Programa e Saída no Console}

O JavaScript não tem uma função \texttt{main}: o arquivo é executado de cima para baixo, e o
código solto depois das classes funciona como programa principal. Comentários de linha começam
com \texttt{//}. Os blocos ficam entre chaves e as instruções terminam com ponto e vírgula. O
\texttt{console.log} imprime um valor e pula uma linha. Textos podem usar aspas simples, aspas
duplas ou crases (\textasciigrave); só com crases o texto vira um \textbf{template literal}, em
que \texttt{\$\{...\}} insere um valor (o equivalente ao \texttt{\#\{...\}} do Ruby).

\begin{lstlisting}[style=javascript]
// comentários de linha começam com //
const banda = 'Linkin Park';   // const: não pode ser reatribuída
let faixas = 2;                // let: pode mudar depois
faixas += 1;
console.log(`A ${banda} tocou ${faixas} músicas.`);  // template literal
console.log('\n--- Fim ---');  // \n = quebra de linha
\end{lstlisting}

\begin{itemize}
    \item \texttt{const} e \texttt{let}: Declaram variáveis. Uma variável \texttt{const} não pode
    receber outro valor depois de criada; uma \texttt{let} pode. As variáveis não têm tipo
    declarado.

    \item \texttt{if (condição) \{ ... \} else \{ ... \}}: Estrutura de decisão. A condição fica
    entre parênteses.

    \item \textbf{Comparações}: \texttt{===} (igual), \texttt{!==} (diferente), \texttt{>},
    \texttt{<}, \texttt{>=} e \texttt{<=}. O \texttt{===} compara o valor e o tipo. Condições
    podem ser combinadas com \texttt{\&\&} (e), \texttt{||} (ou) e \texttt{!} (não).

    \item \texttt{+=} e \texttt{-=}: Somam ou subtraem um valor da variável. Com textos, o
    \texttt{+=} acrescenta o novo texto ao final.
\end{itemize}

\subsection{Declaração de Classe, \texttt{constructor} e \texttt{this}}

Uma classe é declarada com \texttt{class}. O construtor é o método \texttt{constructor},
executado automaticamente quando o objeto é criado com \texttt{new}. Dentro da classe, os
atributos e métodos do próprio objeto são acessados sempre com \texttt{this}:

\begin{lstlisting}[style=javascript]
class Instrumento {
    // roda automaticamente quando o objeto é criado com new
    constructor(nome) {
        this.nome = nome;   // cria o atributo nome no próprio objeto
    }
}

const teclado = new Instrumento('Teclado');  // instanciação com new
console.log(teclado.nome);
\end{lstlisting}

\begin{itemize}
    \item \texttt{constructor}: O construtor, equivalente ao \texttt{initialize} do Ruby e ao
    \texttt{\_\_init\_\_} do Python.

    \item \texttt{this}: Referencia o próprio objeto (como o \texttt{this} do C\#). Diferente do
    Ruby, ele é \textbf{obrigatório}: dentro da classe, um atributo é sempre escrito como
    \texttt{this.nome}.

    \item \texttt{new}: Cria o objeto e chama o \texttt{constructor} com os argumentos
    recebidos.

    \item \textbf{Métodos}: São declarados no corpo da classe só com o nome, os parâmetros e as
    chaves, como em \texttt{tocar() \{ ... \}}, sem nenhuma palavra antes. Na chamada, os
    parênteses são obrigatórios, mesmo sem argumentos: \texttt{guitarra.tocar()}.

    \item \texttt{return}: Devolve um valor para quem chamou o método. Diferente do Ruby, não há
    retorno implícito: sem \texttt{return}, o método não devolve nada.
\end{itemize}

\subsection{Atributos Privados com \texttt{\#} e Getters com \texttt{get}}

Um atributo cujo nome começa com \texttt{\#} é \textbf{privado}: ele só pode ser usado dentro da
classe que o declarou. Diferente do \texttt{\_\_} do Python, que é uma convenção que pode ser
contornada, o JavaScript bloqueia de verdade o acesso. Para que o valor possa ser lido de fora,
a classe cria um \textbf{getter} com \texttt{get}:

\begin{lstlisting}[style=javascript]
class Musico {
    #experiencia = 0;   // privado: declarado no corpo da classe

    constructor(nome, funcao) {
        this.nome = nome;       // público
        this.funcao = funcao;   // público
    }

    get experiencia() {         // getter: usado sem parênteses
        return this.#experiencia;
    }
}

const brad = new Musico('Brad Delson', 'Guitarrista');
console.log(brad.experiencia);    // funciona: usa o getter
console.log(brad.#experiencia);   // SyntaxError: o programa nem roda
\end{lstlisting}

\begin{itemize}
    \item \texttt{\#nome}: Atributo privado. Ele precisa ser \textbf{declarado} no corpo da
    classe, fora do \texttt{constructor}, com ou sem valor inicial (\texttt{\#membros;} ou
    \texttt{\#experiencia = 0;}). Dentro da classe, é usado sempre como \texttt{this.\#nome}, e
    o \texttt{constructor} pode atribuir um valor a ele normalmente.

    \item \textbf{Acesso de fora}: Escrever \texttt{objeto.\#nome} fora da classe é um
    \textbf{erro de sintaxe} (\texttt{SyntaxError}). O JavaScript encontra esse erro antes de
    executar qualquer linha, então o programa inteiro nem começa a rodar.

    \item \textbf{Privado e herança}: O atributo privado pertence só à classe que o declarou. Uma
    classe filha não consegue usar o \texttt{this.\#nome} do pai; para ler o valor, ela usa os
    métodos e getters herdados do pai.

    \item \texttt{get nome() \{ ... \}}: Cria um getter. Ele é declarado como um método, mas é
    usado como um atributo, sem parênteses: \texttt{brad.experiencia}.

    \item \textbf{Somente leitura}: Se a classe tem o getter, mas não tem um setter, uma
    atribuição como \texttt{brad.experiencia = 10}, feita no programa principal, é simplesmente
    ignorada, sem erro, e o valor não muda.
\end{itemize}

Quando o valor privado precisa mudar, a classe oferece um método que faz a alteração de forma
controlada. O exemplo abaixo é complementar e não faz parte do código de referência:

\begin{lstlisting}[style=javascript]
// na classe Musico: altera o atributo privado de forma controlada
ganharExperiencia(pontos) {
    if (pontos > 0) {
        this.#experiencia += pontos;
    }
}
\end{lstlisting}

\subsection{Herança com \texttt{extends} e \texttt{super(...)}}

A classe pai é indicada com \texttt{extends}. O \texttt{constructor} da filha chama
\texttt{super(...)} para executar o \texttt{constructor} do pai, repassando os parâmetros de que
ele precisa (o equivalente ao \texttt{super(...)} do Ruby, ao \texttt{super().\_\_init\_\_(...)}
do Python e ao \texttt{: base(...)} do C\#). Um dos valores repassados pode ser fixo, escrito
diretamente na chamada:

\begin{lstlisting}[style=javascript]
class Vocalista extends Musico {   // Vocalista É UM Musico
    constructor(nome, genero) {
        super(nome, 'Vocalista');  // a função é sempre 'Vocalista'
        this.generoMusical = genero;
    }
}
\end{lstlisting}

Se a filha \textbf{não} define um \texttt{constructor}, ela herda automaticamente o do pai e é
criada com os mesmos parâmetros. O exemplo abaixo é complementar e não faz parte do código de
referência:

\begin{lstlisting}[style=javascript]
class Baixista extends Musico {   // sem constructor próprio
    executarAcao() {
        return `${this.nome} segura a base no baixo!`;
    }
}

// usa o constructor herdado de Musico
const baixista = new Baixista('Phoenix', 'Baixista');
\end{lstlisting}

\begin{itemize}
    \item \texttt{class Filha extends Pai}: Declara a herança (relação ``É UM'').

    \item \texttt{super(...)} \textbf{vem primeiro}: No \texttt{constructor} da filha,
    \texttt{super(...)} deve ser chamado \textbf{antes} de qualquer uso de \texttt{this}. Caso
    contrário, o JavaScript lança um \texttt{ReferenceError}.

    \item \textbf{O que é herdado}: A filha herda os métodos, os getters e os atributos públicos
    do pai. Os atributos privados (\texttt{\#}) continuam acessíveis só pelos métodos do pai.
\end{itemize}

\subsection{Sobrescrita e \texttt{super.metodo()}}

Em JavaScript não existem \texttt{virtual} nem \texttt{override}: basta a classe filha definir
um método com o \textbf{mesmo nome} do método do pai:

\begin{lstlisting}[style=javascript]
// Na Classe Base (Musico):
executarAcao() {
    return `${this.nome} (${this.funcao}) está se preparando...`;
}

// Na Classe Derivada (Vocalista):
executarAcao() {
    return `${this.nome} canta ${this.generoMusical}!`;
}
\end{lstlisting}

Dentro do método sobrescrito, \texttt{super.metodo()} executa a versão do pai (o equivalente ao
\texttt{super} do Ruby, ao \texttt{super().metodo()} do Python e ao \texttt{base.Metodo()} do
C\#). O exemplo abaixo é complementar e não faz parte do código de referência:

\begin{lstlisting}[style=javascript]
// Na Classe Derivada: reaproveita o texto do pai
executarAcao() {
    return super.executarAcao() + ' E solta a voz!';
}
\end{lstlisting}

\subsection{Interfaces Simuladas com \texttt{throw new Error}}

Para simular um contrato, a classe base declara o método obrigatório com um corpo que apenas
lança um erro. Cada filha sobrescreve o método com o seu próprio comportamento:

\begin{lstlisting}[style=javascript]
class Instrumento {
    constructor(nome) {
        this.nome = nome;
    }

    // contrato: toda filha deve sobrescrever tocar()
    tocar() {
        throw new Error(
            `O instrumento ${this.nome} deve implementar o método 'tocar()'`
        );
    }
}

class Bateria extends Instrumento {
    constructor(marca) {
        super('Bateria');
        this.marca = marca;
    }

    tocar() {   // cumpre o contrato
        return `Batendo um ritmo pesado na ${this.marca}!`;
    }
}
\end{lstlisting}

\begin{itemize}
    \item \texttt{throw}: Lança um erro e interrompe o programa. O Node.js exibe a mensagem e a
    linha em que o erro aconteceu.

    \item \texttt{new Error(mensagem)}: Cria o objeto de erro com a mensagem informada.

    \item \textbf{Sem verificação antecipada}: A classe base continua podendo ser instanciada; o
    erro só aparece quando o método é chamado.
\end{itemize}

O exemplo abaixo é complementar e não faz parte do código de referência:

\begin{lstlisting}[style=javascript]
const generico = new Instrumento('Flauta');  // funciona: objeto criado
generico.tocar();   // Error: o erro só aparece aqui, na chamada
\end{lstlisting}

\subsection{Composição e Agregação}

A diferença entre os dois relacionamentos está em \textbf{onde a parte é criada}. Na composição,
o \texttt{new} da parte acontece dentro da classe dona, no \texttt{constructor} ou em qualquer
outro método dela. Na agregação, a parte é criada fora e chega pronta, por parâmetro:

\begin{lstlisting}[style=javascript]
// COMPOSIÇÃO: o Map é criado DENTRO do constructor da Banda
this.#repertorio = new Map();

// AGREGAÇÃO: os músicos são criados FORA e só recebidos pela Banda
const linkinPark = new Banda('Linkin Park',
                             [chester, mike, brad, rob, phoenix]);
\end{lstlisting}

Com uma classe própria, a composição funciona do mesmo jeito: basta criar a parte com
\texttt{new} dentro da classe dona. O exemplo abaixo é complementar e não faz parte do código de
referência:

\begin{lstlisting}[style=javascript]
class Palco {
    #iluminacao;

    constructor(nome) {
        this.nome = nome;
        // composição: a Iluminacao nasce junto com o Palco
        this.#iluminacao = new Iluminacao('Refletores');
    }
}
\end{lstlisting}

\subsection{Arrays: Criação, \texttt{push}, \texttt{forEach} e Arrow Functions}

O \texttt{Array} é a coleção de tamanho dinâmico do JavaScript. Como a tipagem é dinâmica, um
mesmo array pode guardar objetos de classes diferentes, o que possibilita o polimorfismo ao
percorrê-lo:

\begin{lstlisting}[style=javascript]
const membros = [];            // array vazio
membros.push(chester);         // adiciona no final
membros.push(brad);
console.log(membros.length);   // quantidade de elementos

// forEach executa a função para cada elemento
membros.forEach(m => {
    console.log(m.executarAcao());
});

membros.forEach(m => console.log(m.nome));   // versão em uma linha
\end{lstlisting}

\begin{itemize}
    \item \texttt{[]}: Cria um array vazio. \texttt{[a, b]} cria um array já preenchido.

    \item \texttt{array.push(objeto)}: Insere um objeto no final do array (o equivalente ao
    \texttt{<{}<} do Ruby e ao \texttt{append} do Python).

    \item \texttt{array.length}: Retorna a quantidade de elementos.

    \item \texttt{.forEach(função)}: Executa a função uma vez para cada elemento, recebendo o
    elemento da vez como parâmetro.

    \item \textbf{Arrow function}: Forma curta de escrever uma função. Em
    \texttt{m => m.nome}, \texttt{m} é o parâmetro e o valor da expressão depois da seta é
    retornado. Com chaves, como em \texttt{m => \{ ... \}}, o corpo pode ter várias linhas, mas
    só retorna um valor se usar \texttt{return}. Com dois ou mais parâmetros, eles ficam entre
    parênteses: \texttt{(a, b) => ...}.

    \item \texttt{constructor(nome, membrosIniciais = [])}: Um parâmetro pode ter \textbf{valor
    padrão}, usado quando o argumento não é informado.
\end{itemize}

\subsection{Filtrar e Buscar: \texttt{filter}, \texttt{find} e \texttt{instanceof}}

Os arrays também têm métodos que percorrem os elementos e testam uma condição, escrita como
arrow function. O \texttt{find} do exemplo abaixo usa um nome fixo e é complementar; no código
de referência, o nome procurado vem de um parâmetro:

\begin{lstlisting}[style=javascript]
// filter: novo array só com os elementos aprovados pela condição
const guitarristas = membros.filter(m => m.instrumento instanceof Guitarra);

// find: o PRIMEIRO elemento que passa na condição
const musico = membros.find(m => m.nome === 'Brad Delson');
if (musico) {
    console.log(`Encontrado: ${musico.nome}`);
}
\end{lstlisting}

\begin{itemize}
    \item \texttt{.filter(condição)}: Devolve um \textbf{novo} array só com os elementos para os
    quais a condição resultou em \texttt{true}. Se nenhum passar, o array volta vazio.

    \item \texttt{.find(condição)}: Devolve o \textbf{primeiro} elemento que passa na condição.
    Se nenhum passar, devolve \texttt{undefined}, e por isso o resultado costuma ser testado com
    \texttt{if}.

    \item \texttt{objeto instanceof Classe}: Retorna \texttt{true} se o objeto for da classe
    informada ou de uma filha dela (o equivalente ao \texttt{is\_a?} do Ruby e ao
    \texttt{isinstance} do Python).
\end{itemize}

\subsection{Outros Recursos do Código de Referência}

Os recursos abaixo aparecem no código de referência, mas não são necessários nos exercícios:

\begin{itemize}
    \item \texttt{new Map()}: Cria um \textbf{dicionário}, que guarda pares chave e valor.
    \texttt{.set(chave, valor)} adiciona um par, \texttt{.get(chave)} devolve o valor e
    \texttt{.size} informa a quantidade de pares.

    \item \texttt{typeof valor}: Retorna o tipo do valor como texto, por exemplo
    \texttt{'function'}, \texttt{'string'} ou \texttt{'number'}.

    \item \textbf{Objeto literal}: \texttt{\{ nome: 'Teclado', tocar: () => '...' \}} cria um
    objeto na hora, sem classe, com os atributos e métodos escritos entre chaves.

    \item \textbf{Propriedade criada de fora}: Em JavaScript, uma atribuição como
    \texttt{musico.instrumento = ...} cria o atributo no objeto, mesmo que a classe não o
    declare. A linguagem permite, mas isso tira da classe o controle sobre os próprios dados.
\end{itemize}

\section{Código de Referência Completo}

Abaixo está a implementação de referência utilizada para exemplificação do paradigma POO em
JavaScript:

\begin{lstlisting}[style=javascript, caption={Exemplo de Aplicação de POO em JavaScript}]
// CLASSES COMPONENTES E ABSTRATAS (Interfaces)

// simulação de interface/contrato — JavaScript não tem interface nativa
// toda classe que herdar de Instrumento deve implementar tocar()
// o JavaScript não confere isso antes: se a classe não implementar, o
// erro só aparece em tempo de execução, quando tocar() for chamado
class Instrumento {
    constructor(nome) {
        this.nome = nome;
    }

    // VIRTUAL (equivalente)
    // método que as classes filhas vão sobrescrever com o comportamento
    // específico delas
    tocar() {
        throw new Error(
            `O instrumento ${this.nome} deve implementar o método 'tocar()'`
        );
    }
}

// AGREGAÇÃO (Musico -> Instrumento)
// a Guitarra é criada fora (no main) e só depois entregue ao Musico
// pela Banda, por isso o Musico apenas guarda uma referência a ela
// herda de Instrumento e sobrescreve tocar() com seu comportamento
// específico
class Guitarra extends Instrumento {
    constructor(modelo) {
        // super() chama o constructor do pai (Instrumento)
        // equivale ao : base(...) do C#, ao super().__init__() do
        // Python e ao super(...) do Ruby
        super('Guitarra');
        this.modelo = modelo;
    }

    // SOBRESCRITA (override)
    tocar() {
        return `Tocando um riff na ${this.modelo}!`;
    }
}

class Bateria extends Instrumento {
    constructor(marca) {
        super('Bateria');
        this.marca = marca;
    }

    tocar() {
        return `Batendo um ritmo pesado na ${this.marca}!`;
    }
}


// HERANÇA, POLIMORFISMO E ENCAPSULAMENTO

class Musico {
    // ENCAPSULAMENTO
    // # é a sintaxe do JavaScript para atributo privado
    // equivale ao private do C# e do Java: diferente do __ do Python,
    // que pode ser contornado, o JavaScript bloqueia de verdade o acesso
    // fora da classe (escrever musico.#experiencia lá fora é um erro)
    #experiencia = 0;

    constructor(nome, funcao) {
        this.nome = nome;
        this.funcao = funcao;
    }

    // GETTER
    // get é a sintaxe do JavaScript para criar um getter
    // permite ler this.#experiencia de fora com musico.experiencia
    // (sem parênteses)
    get experiencia() {
        return this.#experiencia;
    }

    // VIRTUAL (equivalente)
    executarAcao() {
        return `${this.nome} (${this.funcao}) está se preparando...`;
    }
}

// HERANÇA
// Vocalista herda de Musico com extends
// a relação é: "Vocalista É UM Musico"
class Vocalista extends Musico {
    constructor(nome, genero) {
        super(nome, 'Vocalista');
        this.generoMusical = genero;
    }

    // SOBRESCRITA (override)
    executarAcao() {
        return `${this.nome} canta ${this.generoMusical}!`;
    }
}


// AGREGAÇÃO E COLLECTIONS

class Banda {
    // # tornando os atributos privados
    // membros e repertório não devem ser alterados diretamente de fora
    #membros;
    #repertorio;

    // AGREGAÇÃO
    // membrosIniciais = [] é um parâmetro com valor padrão
    // se nenhum array for passado, começa vazio
    // os Musicos já existiam antes da Banda e continuam existindo se
    // ela sumir
    constructor(nome, membrosIniciais = []) {
        this.nome = nome;
        this.#membros = membrosIniciais;

        // Map é uma estrutura de dados do JavaScript
        // funciona como um dicionário: guarda pares chave -> valor
        // aqui vai guardar título -> duração das músicas
        // COMPOSIÇÃO: o Map é criado aqui dentro e só existe na Banda
        this.#repertorio = new Map();
    }

    adicionarInstrumento(musicoNome, instrumento) {
        // .find() percorre o array e retorna o primeiro elemento que
        // passa na condição
        // m => m.nome === musicoNome é uma arrow function — sintaxe
        // curta do JavaScript
        // equivale a escrever function(m) { return m.nome === musicoNome }
        const musico = this.#membros.find(m => m.nome === musicoNome);
        if (musico) {
            musico.instrumento = instrumento;
            console.log(
                `[SETUP] ${musicoNome} agora tem um(a) ${instrumento.nome}.`
            );
        }
    }

    adicionarMusica(titulo, duracao) {
        // .set() adiciona um par chave -> valor no Map
        this.#repertorio.set(titulo, duracao);
    }

    listarGuitarristas() {
        // .filter() percorre o array e retorna só os elementos que passam
        // na condição
        // instanceof verifica se o objeto é uma instância de determinada
        // classe
        // equivale ao is_a? do Ruby ou isinstance() do Python
        return this.#membros.filter(m => m.instrumento instanceof Guitarra);
    }

    iniciarShow() {
        console.log(`\n--- ${this.nome} INICIA O SHOW! ---`);

        // POLIMORFISMO EM AÇÃO
        // .forEach() percorre o array — equivale ao .each do Ruby ou ao
        // for x in lista do Python
        // o mesmo executarAcao() para todos, mas cada um executa o seu
        // próprio
        this.#membros.forEach(m => {
            let acao = m.executarAcao();

            // typeof verifica o tipo de uma variável em tempo de execução
            // aqui garante que instrumento.tocar realmente é uma função
            // antes de chamar
            if (m.instrumento && typeof m.instrumento.tocar === 'function') {
                acao += ` | ${m.instrumento.tocar()}`;
            }
            console.log(acao);
        });
    }
}


// --- MAIN ---

// criação dos instrumentos: eles existem antes de serem entregues aos
// músicos (agregação)
const guitarraRiff = new Guitarra('PRS Custom');
const guitarraSolo = new Guitarra('Ibanez');
const bateria      = new Bateria('Pearl');

// HERANÇA em ação — cada um é instanciado do seu tipo específico
const chester = new Vocalista('Chester Bennington', 'Nu Metal');
const mike    = new Vocalista('Mike Shinoda',       'Hip Hop');
const brad    = new Musico('Brad Delson', 'Guitarrista');
const rob     = new Musico('Rob Bourdon', 'Baterista');
const phoenix = new Musico('Phoenix',     'Baixista');

// testando getter — acessa #experiencia sem parênteses por causa do get
console.log(`Experiência de Chester: ${chester.experiencia}`);

// AGREGAÇÃO: músicos já existiam antes da banda ser criada
const linkinPark = new Banda('Linkin Park',
                             [chester, mike, brad, rob, phoenix]);

linkinPark.adicionarInstrumento('Brad Delson', guitarraRiff);
linkinPark.adicionarInstrumento('Rob Bourdon', bateria);

// objeto literal passado direto como instrumento — JavaScript permite
// isso
// funciona porque tem o método tocar(), que é tudo que a Banda precisa
// (duck typing)
linkinPark.adicionarInstrumento('Mike Shinoda', {
    nome: 'Teclado',
    tocar: () => 'Adiciona synths marcantes.'
});

linkinPark.adicionarMusica('Somewhere I Belong', 3.5);
linkinPark.adicionarMusica('Fighting Myself',       3.1);

linkinPark.iniciarShow();

console.log('\nGuitarristas Ativos:');
linkinPark.listarGuitarristas().forEach(g => console.log(`- ${g.nome}`));
\end{lstlisting}

\section{Exercícios Práticos do Módulo}

\begin{exercicio}{Exercício 1: Encapsulamento e Herança (Duelo no Ginásio)}
No mundo de \textit{Pokémon}, treinadores levam criaturas chamadas Pokémon para batalhas em
\textbf{ginásios}. Cada Pokémon pertence a uma \textbf{espécie}, tem um \textbf{tipo}, que
define os seus golpes, e uma quantidade de vida. A vida não pode ser alterada diretamente de
fora: só os golpes recebidos em combate podem reduzi-la, e ela nunca fica abaixo de zero.

\textbf{Glossário:}
\begin{itemize}
    \item \textbf{Espécie}: o ``modelo'' de um Pokémon, como Squirtle ou Bulbasaur. No código,
    cada espécie é uma classe.
    \item \textbf{Apelido}: nome que o treinador dá ao seu Pokémon. Dois Squirtles podem ter
    apelidos diferentes.
    \item \textbf{Tipo}: elemento do Pokémon. Todo Squirtle é do tipo Água, e todo Bulbasaur é
    do tipo Planta.
    \item \textbf{Ginásio}: local onde acontecem as batalhas oficiais.
\end{itemize}

Implemente:

\begin{itemize}
    \item \texttt{Pokemon}: o \texttt{constructor} recebe \texttt{apelido}, \texttt{tipo} e
    \texttt{vidaInicial}. \texttt{apelido} e \texttt{tipo} são atributos públicos, e a vida fica
    no atributo privado \texttt{\#vida}, que começa com o valor de \texttt{vidaInicial}. A
    classe tem o getter \texttt{vida} (com \texttt{get}) e nenhum setter: de fora, a vida só pode
    ser lida. Tem também o método \texttt{sofrerDano(pontos)}, que subtrai os pontos da vida sem
    deixá-la negativa (se passar de zero, ela fica em 0) e imprime o apelido e a vida restante; e
    o método \texttt{atacar(alvo)}, que recebe outro Pokémon, imprime um ataque genérico e chama
    \texttt{alvo.sofrerDano(10)}.

    \item \texttt{Squirtle}: herda de \texttt{Pokemon}. O \texttt{constructor} recebe
    \texttt{apelido} e \texttt{vidaInicial} e chama \texttt{super(...)} passando o tipo fixo
    \texttt{'Água'}. Sobrescreve \texttt{atacar(alvo)} para imprimir um golpe de água à sua
    escolha e chamar \texttt{alvo.sofrerDano(20)}.

    \item \texttt{Bulbasaur}: herda de \texttt{Pokemon}, do mesmo jeito, com o tipo fixo
    \texttt{'Planta'}. Sobrescreve \texttt{atacar(alvo)} para imprimir um golpe de planta à sua
    escolha e chamar \texttt{alvo.sofrerDano(15)}.
\end{itemize}

No programa principal, crie um \texttt{Squirtle} e um \texttt{Bulbasaur}, com apelidos à sua
escolha. A vida inicial do Squirtle também é à sua escolha, mas o Bulbasaur deve começar com
30. Faça o Bulbasaur atacar o Squirtle uma vez e o Squirtle atacar o Bulbasaur duas vezes.
Depois, imprima a vida dos dois usando o getter. Por fim, tente imprimir o \texttt{\#vida} de
um dos Pokémon diretamente no programa principal e observe o erro: o programa nem chega a
rodar. Depois, transforme essa linha em comentário (\texttt{//}) para o programa voltar a
rodar.
\end{exercicio}

\begin{exercicio}{Exercício 2: Interface, Composição e Agregação (Esquadrão do Protocolo)}
\textit{Valorant} é um jogo de tiro tático em que dois times de cinco \textbf{agentes} se
enfrentam. Os agentes fazem parte do \textbf{Protocolo VALORANT}, uma organização secreta, e
cada um tem uma \textbf{função} no time, com uma habilidade própria. Todo agente entra na
partida com uma \textbf{pistola}, que é dele e é criada junto com ele (composição). Os agentes
já existem antes de serem reunidos em um \textbf{esquadrão} (agregação).

\textbf{Glossário:}
\begin{itemize}
    \item \textbf{Agente}: personagem jogável de Valorant. Cada agente é conhecido pelo seu
    \textbf{codinome}.
    \item \textbf{Protocolo VALORANT}: organização que reúne os agentes. No código, é a classe
    base de todos eles.
    \item \textbf{Controlador}: função que bloqueia a visão dos inimigos, por exemplo com
    cortinas de fumaça.
    \item \textbf{Iniciador}: função que abre caminho para o time, revelando ou atordoando
    inimigos.
    \item \textbf{Duelista}: função que entra primeiro no combate e busca eliminar os inimigos.
    \item \textbf{Esquadrão}: o time de agentes que joga junto.
\end{itemize}

Implemente:

\begin{itemize}
    \item \texttt{Pistola} (parte composta): o \texttt{constructor} recebe e guarda
    \texttt{modelo}.

    \item \texttt{AgenteProtocolo} (classe base, com contrato e composição): o
    \texttt{constructor} recebe \texttt{codinome} e \texttt{modeloDaPistola}, guarda
    \texttt{codinome} em um atributo público e cria a \texttt{Pistola} \textbf{dentro} dele, com
    o modelo recebido, guardando-a no atributo privado \texttt{\#pistola}. O método
    \texttt{usarHabilidade()} apenas lança um erro com \texttt{throw new Error(...)}, com uma
    mensagem à sua escolha. O método \texttt{entrarEmCombate()} imprime o codinome, o modelo da
    pistola e o texto \textbf{retornado} por \texttt{usarHabilidade()}.

    \item \texttt{Controlador}, \texttt{Iniciador} e \texttt{Duelista}: herdam de
    \texttt{AgenteProtocolo}, \textbf{sem} \texttt{constructor} próprio. Cada uma sobrescreve
    \texttt{usarHabilidade()} para \textbf{retornar} um texto com uma habilidade típica da sua
    função, à sua escolha.

    \item \texttt{Esquadrao} (agregação): o \texttt{constructor} recebe \texttt{nome} e um array
    de agentes já criados, guardado no atributo privado \texttt{\#agentes}. O método
    \texttt{iniciarPartida()} imprime o nome do esquadrão e, com \texttt{forEach}, chama
    \texttt{entrarEmCombate()} de cada agente (polimorfismo). O método
    \texttt{listarControladores()} usa \texttt{filter} com \texttt{instanceof} para
    \textbf{retornar} um array só com os agentes da classe \texttt{Controlador}.
\end{itemize}

No programa principal, crie dois controladores, um iniciador e um duelista, com codinomes e
modelos de pistola à sua escolha. Reúna os quatro em um esquadrão, chame
\texttt{iniciarPartida()} e imprima, com \texttt{forEach}, o codinome de cada agente do array
retornado por \texttt{listarControladores()}. Por fim, crie um \texttt{AgenteProtocolo}
diretamente e, na linha seguinte, chame \texttt{entrarEmCombate()} nele. Observe que o objeto é
criado, mas a chamada gera um erro. Depois, transforme a linha da chamada em comentário
(\texttt{//}) para o programa voltar a rodar.
\end{exercicio}

\begin{exercicio}{Exercício 3: Composição e Agregação (Os Segredos da Cabana)}
Em \textit{Gravity Falls}, dois irmãos gêmeos passam as férias em uma cidadezinha cheia de
mistérios, na casa do tio-avô, que administra a \textbf{Cabana do Mistério}: uma loja de
souvenirs e falsas atrações para turistas. Lá, eles encontram um \textbf{diário misterioso},
escrito por um autor desconhecido, que registra as criaturas sobrenaturais da cidade. Cada
\textbf{ficha} é escrita dentro do diário e só existe nele (composição). Já os
\textbf{funcionários} da Cabana são moradores da cidade, que existiam antes de serem contratados
e continuam existindo fora dela (agregação).

\textbf{Glossário:}
\begin{itemize}
    \item \textbf{Diário misterioso}: caderno que registra as criaturas da cidade. Existem três
    diários, cada um identificado por um número.
    \item \textbf{Ficha de criatura}: registro de uma criatura no diário, com o nome dela e o seu
    nível de perigo.
    \item \textbf{Nível de perigo}: número de 1 (inofensiva) a 5 (muito perigosa). Neste
    exercício, ele não precisa ser validado.
    \item \textbf{Cabana do Mistério}: loja e atração turística da cidade, que tem um dono e
    funcionários.
\end{itemize}

Implemente:

\begin{itemize}
    \item \texttt{FichaDeCriatura}: o \texttt{constructor} recebe \texttt{criatura} (o nome da
    criatura) e \texttt{nivelDePerigo}. O método \texttt{descrever()} \textbf{retorna} um texto
    com os dois dados.

    \item \texttt{DiarioMisterioso}: o \texttt{constructor} recebe \texttt{numero} (o número do
    diário), e a classe tem o atributo privado \texttt{\#fichas}, que começa como um array
    vazio. O método \texttt{anotar(criatura, nivelDePerigo)} cria uma
    \texttt{FichaDeCriatura} \textbf{dentro} do método e a adiciona ao \texttt{\#fichas} com
    \texttt{push}. O método \texttt{exibirFichas()} imprime o número do diário e, com
    \texttt{forEach}, o texto retornado por \texttt{descrever()} de cada ficha.

    \item \texttt{Funcionario}: o \texttt{constructor} recebe \texttt{nome} e \texttt{cargo}.

    \item \texttt{CabanaDoMisterio}: o \texttt{constructor} recebe \texttt{dono} (o nome do
    dono), e a classe tem o atributo privado \texttt{\#funcionarios}, que começa como um array
    vazio. O método \texttt{contratar(funcionario)} adiciona um funcionário já criado ao
    \texttt{\#funcionarios} com \texttt{push}. O método \texttt{abrirAsPortas()} imprime o nome
    do dono e, com \texttt{forEach}, o nome e o cargo de cada funcionário.
\end{itemize}

No programa principal, crie um diário com o número à sua escolha, anote nele duas ou três
criaturas à sua escolha e chame \texttt{exibirFichas()}. Depois, crie dois funcionários
\textbf{antes} da Cabana, à sua escolha. Por fim, crie a Cabana, contrate os dois e chame
\texttt{abrirAsPortas()}. Explique, em um comentário no código, o que é composição e o que é
agregação no seu programa.
\end{exercicio}

\begin{exercicio}{Exercício 4: Revisão Geral (As Aranhas de Yorknew)}
Em \textit{Hunter x Hunter}, \textbf{Hunters} são aventureiros licenciados que dominam o
\textbf{Nen}, a técnica de controlar a própria energia vital, chamada de \textbf{aura}. Todo
usuário de Nen pertence a um \textbf{tipo de Nen}, que define o seu estilo de luta. A
\textbf{Trupe Fantasma}, um grupo de ladrões cujos membros são chamados de \textbf{Aranhas}, foi
avistada em uma cidade, e três Hunters de tipos diferentes vão enfrentá-la.

\textbf{Glossário:}
\begin{itemize}
    \item \textbf{Nen}: técnica de controlar a aura. Todo usuário de Nen sabe atacar com ela.
    \item \textbf{Intensificador}: tipo de Nen que fortalece o próprio corpo, deixando os golpes
    mais fortes.
    \item \textbf{Transmutador}: tipo de Nen que muda as propriedades da aura, por exemplo
    transformando-a em eletricidade.
    \item \textbf{Materializador}: tipo de Nen que cria objetos feitos de aura, como correntes.
    \item \textbf{Aranha}: membro da Trupe Fantasma. Cada Aranha tem um número tatuado, que a
    identifica no grupo.
\end{itemize}

Implemente:

\begin{itemize}
    \item \texttt{UsuarioDeNen} (classe base): o \texttt{constructor} recebe \texttt{nome}. O
    método \texttt{atacar(alvo)} recebe uma \texttt{Aranha} e imprime um ataque genérico de aura
    contra ela, citando o nome dos dois.

    \item \texttt{Intensificador}, \texttt{Transmutador} e \texttt{Materializador}: herdam de
    \texttt{UsuarioDeNen}, \textbf{sem} \texttt{constructor} próprio. Cada uma sobrescreve
    \texttt{atacar(alvo)} para imprimir um ataque típico do seu tipo de Nen, à sua escolha,
    citando o nome dos dois.

    \item \texttt{Aranha}: o \texttt{constructor} recebe \texttt{nome} e \texttt{numero} (o
    número tatuado).

    \item \texttt{TrupeFantasma} (agregação): o \texttt{constructor} recebe \texttt{cidade} e um
    array de Aranhas já criadas, com valor padrão \texttt{[]}, guardado no atributo privado
    \texttt{\#aranhas}. O método \texttt{revelarAranhas()} imprime a cidade e, com
    \texttt{forEach}, o número e o nome de cada Aranha. O método \texttt{buscarAranha(numero)}
    usa \texttt{find} para \textbf{retornar} a Aranha com o número informado.
\end{itemize}

No programa principal, crie duas Aranhas, à sua escolha, \textbf{antes} da Trupe. Crie a Trupe
com uma cidade à sua escolha, chame \texttt{revelarAranhas()} e use
\texttt{buscarAranha(...)} para pegar uma das duas Aranhas. Depois, crie um array com um
usuário de Nen de cada tipo (\texttt{Intensificador}, \texttt{Transmutador} e
\texttt{Materializador}), com nomes à sua escolha, e percorra-o com \texttt{forEach}, fazendo
cada um chamar \texttt{atacar(...)} contra a Aranha encontrada.
\end{exercicio}
